% GENERATED FILE. Edit canonical lessons, then run npm run generate:books.
% canonical-source-hash: fnv1a64:2335a53d056be800
% canonical-lessons: SA-C413-unnayati, SA-C413-sannahyati, SA-C413-goni, SA-C413-prapayati, SA-C413-sampadayati

\chapter{Packing for the Road}
\label{ch:sa-packing-for-the-road}
\hlchaptermodality{413}

\begin{chapteropening}
\textbf{By the end of this chapter:} \emph{I can say what I packed and who will carry it.}

Complete the last lesson of chapter 413: I can say what I packed and who will carry it.
\end{chapteropening}

\section[\texorpdfstring{\sk{उन्नयति} (unnayati)}{उन्नयति (unnayati)}]{\sk{उन्नयति} (unnayati) — raises up; lifts}

\label{lesson:SA-C413-unnayati}



\begin{quote}
Before the new one: say the Sanskrit for passes by, then the Sanskrit for a kos.
\end{quote}

\subsection*{You'll want to know: \sk{उन्नयति}}
\textbf{\sk{उन्नयति}} — \emph{unnayati} — \textquotedblleft{}raises up; lifts\textquotedblright{}.

\begin{quote}
\sk{भारवाहः} \sk{भूमेः} \sk{भारम्} \sk{उन्नयति।}
\end{quote}

\begin{quote}
\emph{bhāravāhaḥ bhūmeḥ bhāram unnayati.}
\end{quote}

\begin{quote}
The porter lifts the load from the ground.
\end{quote}

Past (he …): \sk{सः} \sk{भारम्} \sk{उन्नीतवान्}, \emph{saḥ bhāram unnītavān}.

\begin{cousinweb}
\textbf{\sk{उन्नयति}} (\emph{unnayati}) is built from \emph{ud}, \textquotedblleft{}up\textquotedblright{}, and the root √nī, \textquotedblleft{}to lead\textquotedblright{}, which also gives \emph{nayati}, \textquotedblleft{}leads\textquotedblright{}; MW gives \emph{un-nī} as to lead up and to raise or lift up, with the participle \emph{unnīta} (MW).
\end{cousinweb}

\begin{grammarlens}[title={What you've built}]
One more word for getting around.
\end{grammarlens}

\subsection*{Your turn}
\begin{itemize}
\raggedright
  \item \emph{Say it:} \emph{unnayati}
  \item \emph{Say it:} \emph{unnayati}, once more
  \item \emph{Say it:} \emph{bhāravāhaḥ bhūmeḥ bhāram unnayati.} — The porter lifts the load from the ground.
  \item \emph{Recall:} say the Sanskrit for passes by, then the Sanskrit for a kos, then say \emph{unnayati} again
\end{itemize}

\begin{tcolorbox}[breakable,colback=teal!4,colframe=teal!35!black,title={Before you move on}]
What does \sk{उन्नयति} mean? (\textquotedblleft{}Raises up; lifts\textquotedblright{}.) Say it once more.
\end{tcolorbox}
\section[\texorpdfstring{\sk{सन्नह्यति} (sannahyati)}{सन्नह्यति (sannahyati)}]{\sk{सन्नह्यति} (sannahyati) — gets ready; equips oneself}

\label{lesson:SA-C413-sannahyati}



\begin{quote}
Before the new one: say the Sanskrit for a kos, then the Sanskrit for raises up.
\end{quote}

\subsection*{You'll want to know: \sk{सन्नह्यति}}
\textbf{\sk{सन्नह्यति}} — \emph{sannahyati} — \textquotedblleft{}gets ready; equips oneself\textquotedblright{}.

\begin{quote}
\sk{यात्री} \sk{प्रातः} \sk{यात्रायै} \sk{सन्नह्यति।}
\end{quote}

\begin{quote}
\emph{yātrī prātaḥ yātrāyai sannahyati.}
\end{quote}

\begin{quote}
The traveller gets ready for the journey in the morning.
\end{quote}

Past (he …): \sk{सः} \sk{यात्रायै} \sk{सन्नद्धवान्}, \emph{saḥ yātrāyai sannaddhavān}.

\begin{cousinweb}
\textbf{\sk{सन्नह्यति}} (\emph{sannahyati}) is built from \emph{sam}, \textquotedblleft{}together\textquotedblright{}, and the root √nah, \textquotedblleft{}to bind, to tie\textquotedblright{}: MW gives \emph{saṃ-nah} as to bind together or gird on and to arm oneself or prepare for a task, with the participle \emph{saṃnaddha}, \textquotedblleft{}girt, equipped\textquotedblright{} (MW).
\end{cousinweb}

\begin{grammarlens}[title={What you've built}]
One more word for getting around.
\end{grammarlens}

\subsection*{Your turn}
\begin{itemize}
\raggedright
  \item \emph{Say it:} \emph{sannahyati}
  \item \emph{Say it:} \emph{sannahyati}, once more
  \item \emph{Say it:} \emph{yātrī prātaḥ yātrāyai sannahyati.} — The traveller gets ready for the journey in the morning.
  \item \emph{Recall:} say the Sanskrit for a kos, then the Sanskrit for raises up, then say \emph{sannahyati} again
\end{itemize}

\begin{tcolorbox}[breakable,colback=teal!4,colframe=teal!35!black,title={Before you move on}]
What does \sk{सन्नह्यति} mean? (\textquotedblleft{}Gets ready; equips oneself\textquotedblright{}.) Say it once more.
\end{tcolorbox}
\section[\texorpdfstring{\sk{गोणी} (goṇī)}{गोणी (goṇī)}]{\sk{गोणी} (goṇī) — a sack}

\label{lesson:SA-C413-goni}



\begin{quote}
Before the new one: say the Sanskrit for raises up, then the Sanskrit for gets ready.
\end{quote}

\subsection*{You'll want to know: \sk{गोणी}}
\textbf{\sk{गोणी}} — \emph{goṇī} — \textquotedblleft{}a sack\textquotedblright{}.

\begin{quote}
\sk{भारवाहः} \sk{गोणीं} \sk{पृष्ठे} \sk{वहति।}
\end{quote}

\begin{quote}
\emph{bhāravāhaḥ goṇīṃ pṛṣṭhe vahati.}
\end{quote}

\begin{quote}
The porter carries the sack on his back.
\end{quote}

A feminine noun.

\begin{cousinweb}
\textbf{\sk{गोणी}} (\emph{goṇī}) is the word MW gives for a sack, in Pāṇini's examples and later texts; MW gives no derivation, so its origin is uncertain (MW).
\end{cousinweb}

\begin{grammarlens}[title={What you've built}]
One more word for getting around.
\end{grammarlens}

\subsection*{Your turn}
\begin{itemize}
\raggedright
  \item \emph{Say it:} \emph{goṇī}
  \item \emph{Say it:} \emph{goṇī}, once more
  \item \emph{Say it:} \emph{bhāravāhaḥ goṇīṃ pṛṣṭhe vahati.} — The porter carries the sack on his back.
  \item \emph{Recall:} say the Sanskrit for raises up, then the Sanskrit for gets ready, then say \emph{goṇī} again
\end{itemize}

\begin{tcolorbox}[breakable,colback=teal!4,colframe=teal!35!black,title={Before you move on}]
What does \sk{गोणी} mean? (\textquotedblleft{}A sack\textquotedblright{}.) Say it once more.
\end{tcolorbox}
\section[\texorpdfstring{\sk{प्रापयति} (prāpayati)}{प्रापयति (prāpayati)}]{\sk{प्रापयति} (prāpayati) — delivers; takes to}

\label{lesson:SA-C413-prapayati}



\begin{quote}
Before the new one: say the Sanskrit for gets ready, then the Sanskrit for a sack.
\end{quote}

\subsection*{You'll want to know: \sk{प्रापयति}}
\textbf{\sk{प्रापयति}} — \emph{prāpayati} — \textquotedblleft{}delivers; takes to\textquotedblright{}.

\begin{quote}
\sk{भारवाहः} \sk{गोणीं} \sk{धूमयानं} \sk{प्रापयति।}
\end{quote}

\begin{quote}
\emph{bhāravāhaḥ goṇīṃ dhūmayānaṃ prāpayati.}
\end{quote}

\begin{quote}
The porter takes the sack to the train.
\end{quote}

Past (he …): \sk{सः} \sk{गोणीं} \sk{प्रापितवान्}, \emph{saḥ goṇīṃ prāpitavān}.

\begin{cousinweb}
\textbf{\sk{प्रापयति}} (\emph{prāpayati}) is built from \emph{pra}, \textquotedblleft{}forth\textquotedblright{}, and the causative of the root √āp, \textquotedblleft{}to reach\textquotedblright{}, which gives \emph{prāpnoti}, \textquotedblleft{}reaches\textquotedblright{}; MW gives its participle \emph{prāpita} as conveyed or brought to (MW).
\end{cousinweb}

\begin{grammarlens}[title={What you've built}]
One more word for getting around.
\end{grammarlens}

\subsection*{Your turn}
\begin{itemize}
\raggedright
  \item \emph{Say it:} \emph{prāpayati}
  \item \emph{Say it:} \emph{prāpayati}, once more
  \item \emph{Say it:} \emph{bhāravāhaḥ goṇīṃ dhūmayānaṃ prāpayati.} — The porter takes the sack to the train.
  \item \emph{Recall:} say the Sanskrit for gets ready, then the Sanskrit for a sack, then say \emph{prāpayati} again
\end{itemize}

\begin{tcolorbox}[breakable,colback=teal!4,colframe=teal!35!black,title={Before you move on}]
What does \sk{प्रापयति} mean? (\textquotedblleft{}Delivers; takes to\textquotedblright{}.) Say it once more.
\end{tcolorbox}
\section[\texorpdfstring{\sk{सम्पादयति} (sampādayati)}{सम्पादयति (sampādayati)}]{\sk{सम्पादयति} (sampādayati) — arranges; procures}

\label{lesson:SA-C413-sampadayati}



\begin{quote}
Before the new one: say the Sanskrit for a sack, then the Sanskrit for delivers.
\end{quote}

\subsection*{You'll want to know: \sk{सम्पादयति}}
\textbf{\sk{सम्पादयति}} — \emph{sampādayati} — \textquotedblleft{}arranges; procures\textquotedblright{}.

\begin{quote}
\sk{पिता} \sk{यात्रायै} \sk{यानपत्राणि} \sk{सम्पादयति।}
\end{quote}

\begin{quote}
\emph{pitā yātrāyai yānapatrāṇi sampādayati.}
\end{quote}

\begin{quote}
Father arranges the tickets for the journey.
\end{quote}

Past (he …): \sk{सः} \sk{यानपत्राणि} \sk{सम्पादितवान्}, \emph{saḥ yānapatrāṇi sampāditavān}.

\begin{cousinweb}
\textbf{\sk{सम्पादयति}} (\emph{sampādayati}) is built from \emph{sam}, \textquotedblleft{}together\textquotedblright{}, and the causative of the root √pad, \textquotedblleft{}to go, to fall\textquotedblright{}; MW gives its participle \emph{sampādita} as brought about, accomplished (MW).
\end{cousinweb}

\begin{grammarlens}[title={What you've built}]
One more word for getting around.
\end{grammarlens}

\subsection*{Your turn}
\begin{itemize}
\raggedright
  \item \emph{Say it:} \emph{sampādayati}
  \item \emph{Say it:} \emph{sampādayati}, once more
  \item \emph{Say it:} \emph{pitā yātrāyai yānapatrāṇi sampādayati.} — Father arranges the tickets for the journey.
  \item \emph{Recall:} say the Sanskrit for a sack, then the Sanskrit for delivers, then say \emph{sampādayati} again
\end{itemize}

\begin{tcolorbox}[breakable,colback=teal!4,colframe=teal!35!black,title={Before you move on}]
What does \sk{सम्पादयति} mean? (\textquotedblleft{}Arranges; procures\textquotedblright{}.) Say it once more.
\end{tcolorbox}
